\documentclass[10pt,a4paper]{amsart}
\usepackage[margin=19mm]{geometry}
\usepackage{amsmath,amssymb,mathtools}
\usepackage[scr=rsfs,cal=euler]{mathalfa}
\usepackage{xfrac}
\usepackage[unicode,hidelinks,bookmarks=false]{hyperref}
\newcommand{\ie}{i.e.\ }
\newcommand{\cf}{cf.\ }
\newcommand{\resp}{respectively\ }
\newcommand{\Subsection}{\subsection}
\providecommand{\nobreakdash}{\nobreak\hbox{-}}
\setlength{\emergencystretch}{2em}
\multlinegap0pt
\usepackage{bm}
\usepackage{bbm}
\usepackage{dsfont}
\usepackage{xspace}
\usepackage{cancel}
\usepackage{mathtools}
\usepackage{amsthm}
\makeatletter
\@ifundefined{thm}{\newtheorem{thm}{Theorem}}{}
\makeatother
\newcommand{\C}{\mathbb{C}}
\title[Indeterminate Jacobi operators]{Indeterminate Jacobi operators}
\author{Christian Berg, Ryszard Szwarc}
\date{Source: arXiv:2301.00586v1 (CC BY 4.0)}
\begin{document}
\maketitle

\noindent\emph{Source and licence.} Indeterminate Jacobi operators. Version arXiv:2301.00586v1 (CC BY 4.0), distributed under the Creative Commons Attribution 4.0 International licence, as recorded in the arXiv licence metadata for this exact version and shown on the abstract page https://arxiv.org/abs/2301.00586v1. Full rights notice: licenses/arxiv-2301.00586-CC-BY-4.0.txt.\par\smallskip

\noindent\textit{Editorial scope.} Counted unit: the Thm whose source label is \texttt{thm:deB} in the article (source0.tex lines 793--805, source label \texttt{thm:deB}), delivered together with the article's own complete original proof of it (source0.tex lines 807--839, one \texttt{proof} environment of 33 lines). No mathematics is written, repaired or re-derived: statement and proof are the article's own text line for line. Required context retained uncounted: eq:1 (lines 745--747); eq:3 (lines 756--758); eq:3bis (lines 760--762). Every definition, display and label of those lines is retained unchanged. Omitted from this package and not redistributed: 1--744, 748--755, 759--759, 763--792, 806--806, 840--1242. Nothing inside a retained excerpt is omitted or altered: every retained body is delivered exactly as the article's lines are written, so no omission receipt of any kind accompanies this package. Citations: the retained excerpts carry no citation command at all, so no bibliography entry of the article is transcribed and no reference is redistributed. Reference adaptation: none was needed and none was made; every reference inside the delivered unit resolves inside it, so no unresolved reference or citation is produced. Printed slips kept literal: no printed word, symbol, hypothesis or number of the counted statement or of its proof was altered. Layout and class: the article's own class is \texttt{article} with packages including amssymb, amsmath, amsthm, color, inputenc, srcltx; the wrapper is the lane's pinned \texttt{amsart} editorial wrapper at 10pt on A4, which restates the theorem-like environments the retained text needs and adds the article's own macro definitions that the retained text needs (\texttt{\textbackslash C}), with extra wrapper packages (bm, bbm, dsfont, xspace, cancel, mathtools). The delivered numbering is the unit's own. Assets: no retained span contains a figure environment, a graphics-inclusion command, a PDF-page inclusion, a table or a TikZ picture; no image file is copied into this package, the package declares no asset dependency and no image file is redistributed with it. Licence: version arXiv:2301.00586v1 of this article is distributed under the Creative Commons Attribution 4.0 International licence, as recorded in the arXiv licence metadata for this exact version and shown on the abstract page https://arxiv.org/abs/2301.00586v1. A full rights notice is delivered at licenses/arxiv-2301.00586-CC-BY-4.0.txt. Characters: every retained excerpt is pure ASCII, so the delivered engine, XeLaTeX with its default Latin Modern text font, reports no missing character.
\begin{equation}\label{eq:1}
\frac{p_n(z)-p_n(z_0)}{z-z_0}=\sum_{k=0}^{n-1} a_{n,k}(z_0)p_k(z), \quad z\in\C
\end{equation}

\begin{equation}\label{eq:3}
|a_{n,k}(z_0)|^2\le \left(|p_n(z_0)|^2+|q_n(z_0)|^2\right)\left(|p_k(z_0)|^2+|q_k(z_0)|^2\right).
\end{equation}

\begin{equation}\label{eq:3bis}
\sum_{n=k+1}^\infty |a_{n,k}(z_0)|^2\le (||\mathfrak{p}_{z_0}||^2+||\mathfrak{q}_{z_0}||^2)(|p_k(z_0)|^2+|q_k(z_0)|^2).
\end{equation}

\begin{thm}\label{thm:deB} For $c\in\ell^2$  and $z_0\in\C$ there exists $\xi(c,z_0)\in\ell^2$ such that
\begin{equation}\label{eq:Fcq}
\frac{F_c(z)-F_c(z_0)}{z-z_0}=F_{\xi(c,z_0)}(z) \in\mathcal E, 
\end{equation}
and the coordinates of $\xi(c,z_0)$ are defined by 
\begin{equation}\label{eq:xi}
\xi_k(c,z_0)=\sum_{n=k+1}^\infty c_n a_{n,k}(z_0),\quad k\ge 0.
\end{equation}
Furthermore,
\begin{equation}\label{eq:normxi}
||\xi(c,z_0)|| \le ||c||\left(||\mathfrak{p}_{z_0}||^2+||\mathfrak{q}_{z_0}||^2\right).
\end{equation}
\end{thm}

\begin{proof} The series in \eqref{eq:xi} is absolutely convergent being the product of two $\ell^2$ sequences. Furthermore, by the Cauchy-Schwarz inequality and \eqref{eq:3} we get
$$
|\xi_k(c,z_0)|^2\le ||c||^2 (||\mathfrak{p}_{z_0}||^2+||\mathfrak{q}_{z_0}||^2)\left(|p_k(z_0)|^2+|q_k(z_0)|^2\right),
$$
and therefore $(\xi_k(c,z_0))\in\ell^2$ and \eqref{eq:normxi} holds.

We next find
$$
\frac{F_c(z)-F_c(z_0)}{z-z_0}=\sum_{n=0}^\infty c_n\frac{p_n(z)-p_n(z_0)}{z-z_0},\quad z\neq z_0.
$$ 

 Inserting the expression
\eqref{eq:1} on the right-hand side, we get for $z\neq z_0$
\begin{eqnarray*}
\frac{F_c(z)-F_c(z_0)}{z-z_0}&=&\sum_{n=1}^\infty c_n\sum_{k=0}^{n-1} a_{n,k}(z_0)p_k(z)\\
&=&\sum_{k=0}^\infty p_k(z) \sum_{n=k+1}^\infty c_n a_{n,k}(z_0)\\
&=&\sum_{k=0}^\infty\xi_k(c,z_0) p_k(z),
\end{eqnarray*}
where the rearrangement is possible due to absolute convergence:
\begin{eqnarray*}
\lefteqn{\left|\frac{F_c(z)-F_c(z_0)}{z-z_0}\right| }\\
&\le& \sum_{n=1}^\infty |c_n|\sum_{k=0}^{n-1} |a_{n,k}(z_0)|\,|p_k(z)|\\
&=&\sum_{k=0}^\infty |p_k(z)| \sum_{n=k+1}^\infty |c_n|\,|a_{n,k}(z_0)| \\
&\le& ||c||\sum_{k=0}^\infty |p_k(z)|\left[\sum_{n=k+1}^\infty |a_{n,k}(z_0)|^2\right]^{1/2}\\
&\le& ||c||\sum_{k=0}^\infty |p_k(z)|\left[\left(||\mathfrak{p}_{z_0}||^2+||\mathfrak{q}_{z_0}||^2\right)\left(|p_k(z_0)|^2+|q_k(z_0)|^2\right)\right]^{1/2}\\
&\le& ||c||\left(||\mathfrak{p}_{z_0}||^2+||\mathfrak{q}_{z_0}||^2\right)^{1/2}||\mathfrak{p}_z||\left(\sum_{k=0}^\infty\left(|p_k(z_0)|^2+|q_k(z_0)|^2\right)\right)^{1/2}\\
&=&||c||\,||\mathfrak{p}_z||\left(||\mathfrak{p}_{z_0}||^2+||\mathfrak{q}_{z_0}||^2\right),
\end{eqnarray*}
where we have used \eqref{eq:3bis}. 

It is now clear that the entire functions
$z\mapsto (F_c(z)-F_c(z_0))/(z-z_0)$, with value $F'_c(z_0)$ for $z=z_0$, and $F_{\xi(c,z_0)}(z)$ agree.
\end{proof}

\end{document}
